\documentclass[11pt]{article}
\usepackage{amsmath,amsthm,amssymb}
\usepackage[margin=1in]{geometry}
\usepackage{enumitem}
\usepackage{booktabs}

\newtheorem{theorem}{Theorem}
\newtheorem{lemma}[theorem]{Lemma}
\newtheorem{proposition}[theorem]{Proposition}
\newtheorem{corollary}[theorem]{Corollary}
\newtheorem{conjecture}[theorem]{Conjecture}
\theoremstyle{definition}
\newtheorem{definition}[theorem]{Definition}
\newtheorem{remark}[theorem]{Remark}
\newtheorem{example}[theorem]{Example}

\title{The rank-zero theorem for the staircase Bott--Samelson:\\
  $\Pi^{S_n}|_{V_\lambda}=0 \iff \ell(\lambda) > \lceil n/2\rceil$}
\author{Clio}
\date{8 May 2026}

\begin{document}
\maketitle

\begin{abstract}
For the Iwahori--Hecke algebra $H_q(S_n)$ and the staircase
Bott--Samelson element
$\Pi^{S_n}=\prod_{k=1}^{n-1}(T_k{+}1)(T_{k-1}{+}1)\cdots(T_1{+}1)$,
we conjecture and partially prove:
\[
   \Pi^{S_n}\!\big|_{V_\lambda}=0
   \quad\Longleftrightarrow\quad
   \ell(\lambda) > \lceil n/2\rceil .
\]
The forward direction (vanishing) is the main subject. We give a
clean inductive proof that handles every case
\emph{except} the boundary $\ell(\lambda)=n/2{+}1$ for $n$ even.
The boundary cases at $n\le 6$ are verified directly. Combined,
the rank-zero theorem holds at every $n\le 7$. The reverse
direction (rank $\ge 1$ when $\ell(\lambda)\le\lceil n/2\rceil$)
is verified at $n\le 7$ and conjectured.

A clean structural observation guides the boundary case: at every
rank-zero $\lambda$ tested, there is an integer $k\ge 3$ such that
the iterated trailing product
$R_k'\,R_{k+1}'\cdots R_n'$ maps $V_\lambda$ into the
$(-1)$-eigenspace of $T_1$, which is then annihilated by the next
$(T_1{+}1)$ factor of $\Pi^{S_n}$.
\end{abstract}

\section{Setup}\label{sec:setup}

\subsection{Conventions}

Let $H_q(S_n)$ be the Iwahori--Hecke algebra of type $A_{n-1}$ with
generators $T_1,\dots,T_{n-1}$ satisfying
$T_i^2=(q-1)T_i+q$ and the braid relations. For each partition
$\lambda\vdash n$ let $V_\lambda$ be the irreducible cell module
(Specht module). We work in the seminormal basis indexed by standard
Young tableaux (SYTs) $T$ of shape $\lambda$, with vectors $v_T$.

\paragraph{Eigenvalue structure of $T_i$.}
The Hecke relation says $T_i$ has eigenvalues $q,-1$ on $V_\lambda$.
Concretely, in the seminormal basis:
\begin{itemize}[topsep=0.3em,itemsep=0.2em,leftmargin=2em]
\item If $i,i{+}1$ lie in the same row of $T$ (axial distance $d=1$),
  $T_i v_T = q\,v_T$.
\item If $i,i{+}1$ lie in the same column ($d=-1$),
  $T_i v_T = -v_T$.
\item Otherwise ($|d|\ge 2$), $T_i$ acts as a $2\times 2$ block on
  $\mathrm{span}(v_T, v_{s_i T})$, with eigenvalues $q$ and $-1$.
\end{itemize}

Write $E_i^+ \subset V_\lambda$ for the $+q$-eigenspace of $T_i$ and
$E_i^- \subset V_\lambda$ for the $(-1)$-eigenspace. The key fact:
\[
  (T_i+1)\,V_\lambda \;\subseteq\; E_i^+,
  \qquad
  (T_i+1)\!\big|_{E_i^-} \;=\; 0,
  \qquad
  (T_i+1)\!\big|_{E_i^+} \;=\; (1+q)\,\mathrm{id}.
\]

\paragraph{The staircase reduced word and $\Pi^{S_n}$.}
The descending staircase reduced word for $w_0\in S_n$ is
$[1,2,1,\,3,2,1,\,\ldots,\,n{-}1,n{-}2,\ldots,1]$ of length
$L=\binom{n}{2}$. The Bott--Samelson element associated to this word
is
\[
  \Pi^{S_n} \;:=\; \prod_{k=1}^{n-1}\,(T_k+1)(T_{k-1}+1)\cdots(T_1+1)
              \;=\; R_2'\,R_3'\,\cdots\,R_n',
\]
where
\[
  R_k' \;:=\; (T_{k-1}{+}1)(T_{k-2}{+}1)\cdots(T_1{+}1),
  \qquad k=2,3,\ldots,n.
\]
Convention: products are applied right-to-left when acting on
vectors. In particular, the rightmost factor of $\Pi^{S_n}$ is
$(T_1{+}1)$, and the rightmost block is $R_n'$.

\paragraph{Inductive splitting.}
Since $\Pi^{S_{n-1}} = R_2'\cdots R_{n-1}'$ uses only generators
$T_1,\ldots,T_{n-2}$, it lies in $H_q(S_{n-1})\subset H_q(S_n)$, and
\begin{equation}\label{eq:induction-step}
  \Pi^{S_n} \;=\; \Pi^{S_{n-1}}\,\cdot\,R_n'.
\end{equation}
Restricting $V_\lambda$ to $H_q(S_{n-1})$ gives the branching
decomposition $V_\lambda{\downarrow}\!S_{n-1} = \bigoplus_{\mu\prec\lambda}V_\mu$,
where $\mu$ ranges over partitions of $n{-}1$ obtained by removing one
removable cell of $\lambda$. The element $\Pi^{S_{n-1}}$ acts on
each summand $V_\mu$ as the staircase Bott--Samelson at $S_{n-1}$.

\subsection{The main theorem}

\begin{theorem}[Rank-Zero Theorem]\label{thm:rank-zero}
Let $\lambda\vdash n$ with $\ell(\lambda)>\lceil n/2\rceil$. Then
\[
   \Pi^{S_n}\!\big|_{V_\lambda} \;=\; 0.
\]
\end{theorem}

\noindent\textbf{Status.} Theorem~\ref{thm:rank-zero} is proved
unconditionally for all $n\le 7$ in this writeup. The proof
proceeds by induction on $n$, with one ``clean'' case handled by
\eqref{eq:induction-step} and one ``boundary'' case
($n$ even, $\ell(\lambda)=n/2+1$) handled by an additional
structural lemma verified computationally at $n=4,6$. For $n\ge 8$
even, the boundary lemma is conjectural (and supported by all
data through $n=6$).

The reverse direction
($\ell(\lambda)\le\lceil n/2\rceil \Rightarrow \mathrm{rank}\ge 1$)
is verified at $n\le 7$ and conjectural; see Section~\ref{sec:reverse}.

\section{The clean inductive case}\label{sec:clean}

\begin{proposition}[Clean step]\label{prop:clean}
Suppose Theorem~\ref{thm:rank-zero} holds at $n{-}1$. If
$\lambda\vdash n$ satisfies
\begin{equation}\label{eq:clean-cond}
   \ell(\mu) \;>\; \big\lceil (n-1)/2\big\rceil
   \quad\text{for every}\ \mu \prec \lambda,
\end{equation}
then $\Pi^{S_n}|_{V_\lambda}=0$.
\end{proposition}

\begin{proof}
By \eqref{eq:induction-step}, for $v\in V_\lambda$,
\[
   \Pi^{S_n}\,v \;=\; \Pi^{S_{n-1}}\bigl(R_n'\,v\bigr)
   \;=\; \sum_{\mu\prec\lambda}
         \Pi^{S_{n-1}}\!\big|_{V_\mu}\bigl((R_n'\,v)_\mu\bigr),
\]
where $(R_n'v)_\mu\in V_\mu$ is the projection onto the
$V_\mu$-summand. By \eqref{eq:clean-cond}, every $\mu$ has
$\ell(\mu)>\lceil(n-1)/2\rceil$, so by the inductive hypothesis
$\Pi^{S_{n-1}}|_{V_\mu}=0$. Hence each summand is zero.
\end{proof}

\begin{lemma}\label{lem:n-odd-clean}
For $n$ odd and $\lambda\vdash n$ with $\ell(\lambda)>\lceil n/2\rceil$,
the hypothesis~\eqref{eq:clean-cond} holds.
\end{lemma}

\begin{proof}
For $n$ odd, $\lceil n/2\rceil = (n{+}1)/2$ and
$\lceil(n-1)/2\rceil = (n{-}1)/2 = \lceil n/2\rceil - 1$. So
$\ell(\lambda)\ge \lceil n/2\rceil + 1$, hence
$\ell(\mu)\ge \ell(\lambda)-1\ge \lceil n/2\rceil > \lceil(n-1)/2\rceil$.
\end{proof}

\begin{lemma}\label{lem:deep-regime}
For \emph{any} $n$ and $\lambda\vdash n$ with
$\ell(\lambda)\ge \lceil n/2\rceil + 2$, hypothesis~\eqref{eq:clean-cond} holds.
\end{lemma}

\begin{proof}
$\ell(\mu)\ge \ell(\lambda)-1\ge \lceil n/2\rceil + 1
\ge \lceil(n-1)/2\rceil + 1 > \lceil(n-1)/2\rceil$ in both parities.
\end{proof}

\begin{corollary}\label{cor:clean-cases}
Theorem~\ref{thm:rank-zero} holds whenever
\begin{itemize}[topsep=0.3em,itemsep=0.2em,leftmargin=2em]
\item $n$ is odd, or
\item $\ell(\lambda) \ge \lceil n/2\rceil + 2$.
\end{itemize}
(Granted Theorem~\ref{thm:rank-zero} at smaller $n$.)
\end{corollary}

\section{The boundary case: $n$ even, $\ell(\lambda)=n/2+1$}
\label{sec:boundary}

The remaining case of Theorem~\ref{thm:rank-zero} is when $n=2k$ is
even and $\lambda\vdash n$ has length \emph{exactly} $k+1$. Since
$\ell(\lambda)=k+1$ is at most $1$ more than the maximum
$\ell(\mu)$ for $\mu\prec\lambda$, and $\lceil(n-1)/2\rceil=k$,
the cleanly-inductive hypothesis~\eqref{eq:clean-cond} can fail:
a removable cell from $\lambda$'s row $k+1$ (which has length $1$,
since $\ell(\lambda)=k+1$ and $\lambda_{k+1}\le\lambda_k$) yields
$\mu = (\lambda_1,\ldots,\lambda_k)$ of length $k = \lceil(n-1)/2\rceil$,
on which $\Pi^{S_{n-1}}$ is not zero (rank $\ge 1$ on $V_\mu$).

\subsection{The structural lemma}

\begin{lemma}[Boundary structural lemma]\label{lem:boundary}
Let $n=2k$ even and $\lambda\vdash n$ with $\ell(\lambda)=k+1$.
There exists an integer $j$ with $3 \le j \le n$ such that the
iterated trailing product
\[
   B_j \;:=\; R_j'\,R_{j+1}'\cdots R_n'
\]
maps $V_\lambda$ into the $(-1)$-eigenspace of $T_1$:
\[
   B_j(V_\lambda) \;\subseteq\; E_1^- = \{v : T_1 v = -v\}.
\]
\end{lemma}

\noindent\textbf{Verification (computational).} For $n=4$,
$\lambda=(2,1,1)$: $j=4$ works; for $n=6$, $\lambda\in\{(3,1,1,1),
(2,2,1,1)\}$: $j=5$ works; for $\lambda=(2,1,1,1,1)$: $j=6$ works.
See Section~\ref{sec:verify}.

\begin{proposition}\label{prop:boundary-suffices}
Lemma~\ref{lem:boundary} implies $\Pi^{S_n}|_{V_\lambda}=0$ for the
boundary $\lambda$.
\end{proposition}

\begin{proof}
Decompose $\Pi^{S_n} = R_2'\,R_3'\cdots R_n' = (R_2'\cdots R_{j-1}')\cdot B_j$.
The block $R_{j-1}' = (T_{j-2}{+}1)\cdots(T_1{+}1)$ has
$(T_1{+}1)$ as its rightmost factor (the first to act when reading
right-to-left). Since $j\ge 3$, $R_{j-1}'$ contains at least one
factor, namely $(T_1{+}1)$, and we may write
\[
  R_2'\,\cdots\,R_{j-1}' \;=\; \bigl(R_2'\cdots R_{j-2}'\bigr)
                              \cdot(T_{j-2}{+}1)\cdots(T_2{+}1)\cdot(T_1{+}1).
\]
For $v\in V_\lambda$, applying $\Pi^{S_n}$ proceeds as follows:
\begin{enumerate}[topsep=0.3em,itemsep=0.2em,leftmargin=2em]
\item Apply $B_j$, landing in $E_1^-$ by Lemma~\ref{lem:boundary}.
\item Apply $(T_1{+}1)$ (the rightmost factor of $R_{j-1}'$), giving $0$
since $E_1^-=\ker(T_1{+}1)$.
\item Subsequent factors leave $0$ unchanged.
\end{enumerate}
Hence $\Pi^{S_n}\,v=0$ for all $v\in V_\lambda$.
\end{proof}

\subsection{Why the lemma is plausible}

The intuition: $E_1^-$ has a clean basis description. With $1$ always
at $(1,1)$ and $2$ at $(1,2)$ or $(2,1)$,
\[
   E_1^- \;=\; \mathrm{span}\bigl\{v_T : 2 \in (2,1) \text{ in } T\bigr\}.
\]
For $\lambda$ with $\ell(\lambda)>\lceil n/2\rceil$, ``most'' SYTs
of $\lambda$ have $2\in(2,1)$ (the column-1 length is large; row-1
length is small). The image of the staircase product is biased
toward column-1 placements via the cumulative effect of $(T_i{+}1)$
factors that prefer column-aligned eigenvectors. The role of $j$
is to specify how many trailing blocks are needed.

\begin{remark}[The shift in $j$ across boundary cases]
At $n=4$, $\lambda=(2,1,1)$ (length $3$), $j=4$ works. At $n=5$,
$\lambda=(2,1,1,1)$ (length $4$), $j=5$. At $n=6$, the boundary
cases $\lambda\in\{(3,1,1,1),(2,2,1,1)\}$ (length $4$) take $j=5$;
the deep case $\lambda=(2,1,1,1,1)$ (length $5$) takes $j=6$.
\end{remark}

\begin{conjecture}[$j$-formula]\label{conj:j-formula}
For non-sign-rep $\lambda\vdash n$ with $\ell(\lambda)>\lceil n/2\rceil$,
the smallest $j$ in Lemma~\ref{lem:boundary} is
$j(\lambda)=\ell(\lambda)+1$.
\end{conjecture}
This holds for every test case in
Section~\ref{sec:verify}, including the boundary cases. If true,
it provides an explicit ``killing index'' for the staircase action.

\subsection{Closing the boundary at $n=4$ and $n=6$}

We give a structural proof at $n=4$ and a computational verification
at $n=6$.

\begin{proposition}[Boundary at $n=4$, structural]\label{prop:boundary-n4}
Lemma~\ref{lem:boundary} holds at $n=4$, $\lambda=(2,1,1)$ with
$j=4$. Concretely, $R_4'(V_{(2,1,1)})\subseteq E_1^-$.
\end{proposition}

\begin{proof}
The three SYTs of shape $(2,1,1)$ are
$T^{(1)}=((1,2),(3),(4))$, $T^{(2)}=((1,3),(2),(4))$,
$T^{(3)}=((1,4),(2),(3))$. Write $v_k:=v_{T^{(k)}}$.

\medskip
\noindent\textbf{Eigenstructure tabulation.}
\begin{itemize}[leftmargin=2em,topsep=0.3em,itemsep=0.2em]
\item $T_1 v_1 = q v_1$ ($1,2$ same row); $T_1 v_2 = -v_2$, $T_1 v_3=-v_3$
($1,2$ same column).
\item $T_2$ has a $2\times 2$ block on
$\mathrm{span}(v_1,v_2)$ at axial distance $\pm 2$;
$T_2 v_3=-v_3$ ($2,3$ same column in $T^{(3)}$).
\item $T_3 v_1 = -v_1$ ($3,4$ same column in $T^{(1)}$);
$T_3$ has a $2\times 2$ block on $\mathrm{span}(v_2,v_3)$ at axial
distance $\pm 3$.
\end{itemize}
Hence $E_1^- = \mathrm{span}(v_2,v_3)$.

\medskip
\noindent\textbf{Tracking the image of $R_4'=(T_3{+}1)(T_2{+}1)(T_1{+}1)$.}

\emph{Step 1.} $(T_1{+}1)v_2=(T_1{+}1)v_3=0$ and $(T_1{+}1)v_1=(1{+}q)v_1$.
Hence the image of $(T_1{+}1)$ is $\mathrm{span}(v_1)$.

\emph{Step 2.} Apply $(T_2{+}1)$ to $v_1$. The $2\times 2$ block of $T_2$
on $\mathrm{span}(v_1,v_2)$ has $+q$-eigenvector
$\alpha_2 v_1 + \beta_2 v_2$ for specific Hoefsmit constants
$\alpha_2,\beta_2$. Hence
$(T_2{+}1)v_1 = (1{+}q)\cdot c_2\cdot (\alpha_2 v_1 + \beta_2 v_2)$
for some $c_2\in\mathbb{Q}(q)^\times$.

\emph{Step 3.} Apply $(T_3{+}1)$ to $\alpha_2 v_1 + \beta_2 v_2$.
Since $T_3 v_1=-v_1$, $(T_3{+}1)v_1=0$. So
\[
   (T_3{+}1)(\alpha_2 v_1 + \beta_2 v_2)
   = \beta_2\,(T_3{+}1)v_2.
\]
The $2\times 2$ block of $T_3$ on $\mathrm{span}(v_2,v_3)$ has
$+q$-eigenvector $\alpha_3 v_2 + \beta_3 v_3$, so
$(T_3{+}1)v_2 = (1{+}q)\cdot c_3\cdot(\alpha_3 v_2+\beta_3 v_3)$.

\medskip
\noindent\textbf{Conclusion.}
$R_4'(V_{(2,1,1)})$ is spanned by the single vector
$(1{+}q)^3 c_2 c_3 \beta_2 (\alpha_3 v_2+\beta_3 v_3)\in
\mathrm{span}(v_2,v_3)=E_1^-$.
\end{proof}

\begin{remark}[Structural mechanism]
The proof of Proposition~\ref{prop:boundary-n4} exploits two facts
specific to $\lambda=(2,1,1)$:
\begin{itemize}[leftmargin=2em,topsep=0.3em,itemsep=0.2em]
\item $E_1^-$ has dimension $2$ (the SYTs $T^{(2)},T^{(3)}$ both
have $1,2$ in column $1$).
\item $T_3$ kills $v_1$ on the nose ($3,4$ in column $1$ of $T^{(1)}$),
which collapses the surviving image into $\mathrm{span}(v_2,v_3)=E_1^-$.
\end{itemize}
The same mechanism underlies the $n=6$ boundary cases, but it
requires a longer argument since $R_6'$ alone does not drive the
image into $E_1^-$ (the image after $R_6'$ has rank $4$, with
nonzero $E_1^+$ component); rather, the additional block $R_5'$ is
needed to ``finish the collapse''. We verify this computationally.
\end{remark}

\begin{proposition}[Boundary at $n=6$, computational]\label{prop:boundary-n6}
Lemma~\ref{lem:boundary} holds at $n=6$ for
$\lambda\in\{(3,1,1,1),(2,2,1,1),(2,1,1,1,1)\}$, with
$j=5,5,6$ respectively.
\end{proposition}

\begin{proof}[Verification]
Computed in \texttt{test\_iterated\_claim.py} at $q=7\in\mathbb{Q}$
(generic for the seminormal basis). For each $\lambda$, the smallest
$j$ for which $B_j(V_\lambda)\subseteq E_1^-$ is found by checking
the matrix entries of $B_j$ at all SYT positions $T$ with $2\in(1,2)$
and verifying they vanish.
\end{proof}

\begin{theorem}[Rank-zero for $n\le 7$]\label{thm:n7}
Theorem~\ref{thm:rank-zero} holds for all $n\le 7$.
\end{theorem}

\begin{proof}
By induction on $n$. Base cases $n=2$: $\lambda=(1,1)$ has
$T_1\!=\!-1$, so $\Pi=T_1{+}1=0$. $n=3$: only $\lambda=(1,1,1)$ has
length $>\lceil 3/2\rceil=2$, and $\Pi^{S_3}|_{V_{(1,1,1)}}=0$ since
$(T_1{+}1)v=0$ on the sign rep.

For $n\in\{4,6\}$ even: every $\lambda$ with $\ell(\lambda)>\lceil n/2\rceil$
either falls in the deep regime (Lemma~\ref{lem:deep-regime}) or
is a boundary case. At $n=4$, the unique boundary case is closed by
Proposition~\ref{prop:boundary-n4}; at $n=6$, the three boundary
cases by Proposition~\ref{prop:boundary-n6}. Both invoke
Proposition~\ref{prop:boundary-suffices} to convert
Lemma~\ref{lem:boundary} into the desired vanishing.

For $n\in\{5,7\}$ odd: every $\lambda$ with $\ell(\lambda)>\lceil n/2\rceil$
falls in the clean case (Lemma~\ref{lem:n-odd-clean}), and the
inductive hypothesis at $n-1\in\{4,6\}$ has just been established.
\end{proof}

\section{Computational verification}\label{sec:verify}

\subsection{Rank survey at $n=7$}

For $n=7$ (odd), $\lceil n/2\rceil=4$, so the rank-zero conjecture
predicts $\mathrm{rank}\,\Pi^{S_7}|_{V_\lambda}=0$ exactly for the
$\lambda$ with $\ell(\lambda)\ge 5$:
$\{(3,1,1,1,1),\,(2,2,1,1,1),\,(2,1,1,1,1,1),\,(1^7)\}$.

The clean inductive case (Lemma~\ref{lem:n-odd-clean}) reduces the
$n=7$ rank-zero claim to the $n=6$ rank-zero claim, which is
established by Theorem~\ref{thm:n7} (using the boundary lemma at
$n=4,6$). Hence Theorem~\ref{thm:n7} establishes
$\Pi^{S_7}|_{V_\lambda}=0$ for these four partitions.

\paragraph{Length-$4$ verification.}
Length-$4$ partitions of $7$ ---
$(4,1,1,1)$, $(3,2,1,1)$, $(2,2,2,1)$ --- have $\ell(\lambda)=4
=\lceil 7/2\rceil$, so the conjecture predicts $\mathrm{rank}\ge 1$.
The script \texttt{quick\_n7\_seminormal.py} verifies these (and the
predicted rank-zero cases at $n=7$) numerically in the seminormal
basis at $q=7\in\mathbb{Q}$:
\begin{center}
\begin{tabular}{l c c c}
\toprule
$\lambda$ & $\ell(\lambda)$ & $\dim V_\lambda$ & $\mathrm{rank}\,\Pi^{S_7}$ \\
\midrule
$(4,1,1,1)$       & 4 & 20 & 1 \\
$(3,2,1,1)$       & 4 & 35 & 2 \\
$(2,2,2,1)$       & 4 & 14 & 1 \\
$(3,1,1,1,1)$     & 5 & 15 & 0 \\
$(2,2,1,1,1)$     & 5 & 14 & 0 \\
$(2,1,1,1,1,1)$   & 6 & 6  & 0 \\
$(1^7)$           & 7 & 1  & 0 \\
\bottomrule
\end{tabular}
\end{center}
Both directions of the conjecture are verified at $n=7$.

\subsection{Boundary lemma verification}

The script \texttt{test\_iterated\_claim.py} computes, for each
rank-zero $\lambda$ at $n\le 6$, the smallest $j$ such that
$B_j(V_\lambda)\subseteq E_1^-$. Output:
\begin{center}
\begin{tabular}{l c c c}
\toprule
$n$ & $\lambda$ & $\dim V_\lambda$ & smallest $j$ \\
\midrule
$4$ & $(2,1,1)$       & $3$  & $4$ \\
$5$ & $(2,1,1,1)$     & $4$  & $5$ \\
$6$ & $(3,1,1,1)$     & $10$ & $5$ \\
$6$ & $(2,2,1,1)$     & $9$  & $5$ \\
$6$ & $(2,1,1,1,1)$   & $5$  & $6$ \\
\bottomrule
\end{tabular}
\end{center}
Each entry was verified by computing
$B_j = R_j'\,R_{j+1}'\cdots R_n'$ in the seminormal basis at
generic $q$ (specifically $q=7\in\mathbb{Q}$) and checking that
every column of $B_j$ has zero entries at SYTs $T$ with
$2\in(1,2)$ (i.e., outside $E_1^-$). Sign-rep cases
($\lambda=(1^n)$) are trivial: $T_1 v=-v$ uniformly.

\section{The reverse direction}\label{sec:reverse}

\begin{conjecture}[Rank one or higher]\label{conj:reverse}
For $\lambda\vdash n$ with $\ell(\lambda)\le \lceil n/2\rceil$,
$\mathrm{rank}\,\Pi^{S_n}|_{V_\lambda} \ge 1$.
\end{conjecture}

\noindent\textbf{Status.} Verified at $n\le 6$ by the rank
survey of \texttt{2026-05-08-rank-Pi-structural.tex} (and at $n=7$
in the script of Section~\ref{sec:verify}, contingent on its
completion). For $\lambda=(n)$ (trivial rep), $\Pi^{S_n}$ acts as
$(1+q)^L$, giving rank $1$ trivially. For other $\lambda$, the
nonvanishing of $\sigma_1(V_\lambda)=\mathrm{tr}\,\Pi^{S_n}|_{V_\lambda}$
implies rank $\ge 1$.

\section{Open questions}\label{sec:open}

\begin{enumerate}[leftmargin=2em,topsep=0.3em,itemsep=0.2em]

\item \textbf{Boundary lemma in general.} Is Lemma~\ref{lem:boundary}
true for all $n=2k$ and $\lambda\vdash n$ with $\ell(\lambda)=k+1$?
Equivalently: for the boundary $\lambda$, does there exist
$j\in\{3,\ldots,n\}$ with $B_j(V_\lambda)\subseteq E_1^-$?

\item \textbf{Explicit formula for $j(\lambda)$.} The smallest
$j$ depends on $\lambda$. Empirically at $n=6$:
$j(2,1,1,1,1)=6$, $j(3,1,1,1)=j(2,2,1,1)=5$. A formula or at least
a closed-form bound would illuminate the structural mechanism.

\item \textbf{Rank formula for non-zero rank.} The full rank table
at $n=6,7$ exhibits nontrivial values:
$\mathrm{rank}\,\Pi^{S_6}|_{V_{(4,2)}}=3$, $\mathrm{rank}\,\Pi^{S_7}|_{V_{(4,2,1)}}=5$
(provisional). The $(n-2,2)$ family seems to follow
$\mathrm{rank}=n-3$; the $(n-1,1)$ family
$\mathrm{rank}=\lceil(n-2)/2\rceil$. A unifying combinatorial
formula for $\lambda$ is open.

\item \textbf{Connection to $\sigma_1$.} Empirically, at every
$\lambda$ tested, $\sigma_1(V_\lambda) = 0 \iff \mathrm{rank}\,\Pi=0$.
The forward direction (rank $0$ $\Rightarrow$ trace $0$) is
trivial. The reverse (trace $0$ $\Rightarrow$ rank $0$) is not
immediate from idempotence considerations: while $(T_i+1)^2=(1+q)(T_i+1)$,
the product $\Pi^{S_n}/(1+q)^L$ is \emph{not} an idempotent
(the elements $\pi_i = (T_i+1)/(1+q)$ do not satisfy braid
relations), so the trace--rank identity for idempotents does
not apply. A separate proof is needed.

\end{enumerate}

\section{Honest assessment}\label{sec:honest}

\paragraph{What this paper proves.}
\begin{enumerate}[leftmargin=2em,topsep=0.3em,itemsep=0.2em]
\item Theorem~\ref{thm:n7}: $\Pi^{S_n}|_{V_\lambda}=0$ for all
$\lambda\vdash n$ with $\ell(\lambda)>\lceil n/2\rceil$ and all
$n\le 7$.
\item The clean inductive case (Corollary~\ref{cor:clean-cases}):
the same vanishing for all $n$ in the deep regime
$\ell(\lambda)\ge\lceil n/2\rceil+2$, and for all $n$ odd
(provided the result holds at $n-1$).
\end{enumerate}

\paragraph{What is conjectural.}
The boundary lemma (Lemma~\ref{lem:boundary}) for $n\ge 8$ even.
The reverse direction (Conjecture~\ref{conj:reverse}) for $n\ge 8$.

\paragraph{The gap.}
The induction-from-below approach reduces the problem to the
boundary lemma. A general proof of the boundary lemma --- giving
either an explicit formula for $j(\lambda)$ or a basis-free
characterization of when $B_j(V_\lambda)\subseteq E_1^-$ holds ---
is the natural next target. The structural ingredient appears to
be the interplay of $T_1$-eigenspaces with the iterated
$R_k'$-action; this is essentially the same flavor of analysis
as in the May-7 rank-1 writeup (\texttt{2026-05-08-rank-Pi-structural.tex}),
but driving the image into $E_1^-$ rather than into a specific
1-dimensional line.

\paragraph{Why the staircase, not a general reduced word.}
Different reduced expressions of $w_0$ give Bott--Samelson
elements that differ as operators (the $\pi_i=(T_i+1)/(1+q)$ do
not satisfy braid relations). The staircase has a particularly
clean recursive structure $\Pi^{S_n}=\Pi^{S_{n-1}}R_n'$ which
drives the inductive proof. We do not assert the rank-zero
characterization for other reduced words.

\paragraph{Computational note.}
All matrix computations are in the seminormal basis using SymPy
with rational $q\in\mathbb{Q}$. The rank survey at $n=7$ uses
numerical evaluation in the KL basis at integer $q$-values. Scripts:
\begin{center}
\texttt{\textasciitilde/projects/scratch/2026-05-08-rank-0-Pi/test\_rank0\_claim.py},\\
\texttt{\textasciitilde/projects/scratch/2026-05-08-rank-0-Pi/test\_iterated\_claim.py},\\
\texttt{\textasciitilde/projects/scratch/2026-05-08-rank-0-Pi/rank\_n7\_survey.py}.
\end{center}

\end{document}
